\label{ch:lit_review}



\section{Классические модели с отказами и свойство инвариантности}
\label{sec:insensitivity_general}

В теории массового обслуживания(далее ТМО) наиболее сложными для анализа моделями являются самые общие модели. Чем больше факторов влияет на работу системы массового обслуживания (далее СМО) и чем разнообразнее формы распределения, тем сложнее аналитически анализировать такие системы. А.К. Эрланг первым предложил использовать марковские случайные процессы для описания СМО\cite{bocharov_pechinkin_tmo,adan_resing2015_queueing_systems}. 

Самый простой такой процесс с пуассоновским входящим потоком заявок и экспоненциальным временем обслуживания и бесконечной очередью описывается в нотации Кендалла как $M/M/1$\cite{bocharov_pechinkin_tmo,adan_resing2015_queueing_systems}. Первые 2 буквы обозначают марковость входящего потока и времени обслуживания, задавая фундаментальное свойство независимости каждого следующего события от прошлого, а единица обозначает, что 1 прибор обслуживает 1 заявку. В таких системах можно исследовать стационарное распределение вероятностей состояний, среднее число заявок в системе и среднюю длину очереди, среднее время пребывания заявки в системе и среднее время ожидания, пропускную способность системы. 

Добавляя в модель отсутствие очереди и потери заявок усложняет систему и в нотации она выглядит как $M/M/n/0$\cite{bocharov_pechinkin_tmo,adan_resing2015_queueing_systems}, где n - количество независимых приборов, а 0 --- очередь, точнее ее отсутствие. Также используется нотация $M/M/n/n$ аналогичная по смыслу, где $n$ на четвертом месте обозначает число заявок в системе, после которого наступает отказ. В случае, когда заявки могут теряться появляется возможность исследовать вероятность отказа, чувствительность к числу приборов. По прежнему входящий поток является пуассоновским, то есть вероятность поступления одной заявки на малом интервале $\delta$ равна $\delta\lambda + o(\delta)$, где $\lambda$ --- интенсивность потока, а вероятность поступления двух и более заявок имеет порядок малости $o(\delta)$. А время обслуживания каждой заявки заявки на каждом из приборов --- случайная величина, распределенная по экспоненциальному закону с параметром $\mu$. 

В реальных телекоммуникационных и вычислительных системах, такое строгое ограничение на входящий поток и время работы зачастую не работает. Если взять более широкую систему $M/G/n/n$, где время работы имеет произвольное распределение $G$ (General) с конечным средним $b = \mathbb{E}(S)$, а в экспоненциальном случае $\mathbb{E}S = 1/\mu$, так что у системы пропадает свойство марковости. Однако, такая система, как было доказано Б.А. Севастьяновым\cite{sevastyanov1957_ergodic_refusals,burman1981_insensitivity}, обладает свойством инвариантности, то есть может быть сведена к $M/M/n/n$, а вероятность отказа зависит исключительно от первого момента времени обслуживания\cite{sevastyanov1957_ergodic_refusals,burman1981_insensitivity} интернет трафика. 

Свойство инвариантности также распространяется на сиситемы типа $M/G/1-PS$\cite{burman1981_insensitivity,adan_resing2015_queueing_systems} (разделение вычислительной мощности процессора между $k$ заявками), $M/G/\infty$ и системы с инверсионным порядком обслуживания и вытеснением (LCFS-PR, Last-Come First-Served with Preemptive Resume), но не распространяется на классическую модель $M/G/1/\infty$ или на непуассоновский входящий поток $G/G/n/n$. Данный факт обуславливает выбор СМО с потерями для данной работы, в ходе работы для пуассоновсокого входящего потока будет эмпирически проверено свойство инвариантности. 
%Начинаем с базы. Вводим понятие многоканальной системы с отказами (M/G/N/N). Рассказываем про формулу Эрланга. Затем вводим свойство инвариантности (нечувствительности): почему стационарное распределение зависит только от среднего времени обслуживания, но не от его дисперсии.

% Блок 3: границы применимости известных результатов.

% Цитировать Бочаров и Печинкин (учебник), Burman (Insensitivity in Queueing Systems), Brumelle (A Generalization of Erlang’s Loss System), Adan & Resing

\section{Математические модели систем с ограниченными ресурсами}
\label{sec:limited_resources}

Ранее проанализированные системы типа $M/G/n/n$ также не учитывают возможность наличия у заявки требования помимо канала (или нескольких каналов обслуживания), но и требования к ресурсу. Так, например, у современных сетей 5G для среза URLLC заявка может иметь требование $n$ к количеству виртуальных процессоров и к ресурсу $R$ --- объем оперативной памяти. Нехватка любого из требуемых ресурсов может стать причиной потери заявки. Далее нотация Кендалла становится неполной для иллюстрации, например, распределения ресурсного требования и становится проще задать систему описанием набора объектов: число серверов, число типов ресурсов, общий ресурс, распределения межприходных интервалов, обслуживания и ресурсных требований.

Для удобного представления подобных моделей как правило придерживаются идеи о том, что входящая заявка окупирует один прибор и имеет требование к ресурсу (число или вектор для мультиресурсной системы). Тогда система состоит из $N$ обслуживающих приборов (серверов) и $M$ типов ресурсов. Общий доступный объем ресурсов в системе задается вектором $R=(R_{1},R_{2},...,R_{M})$, а входящий поток может состоять из $L$ различных типов заявок. Времена между поступлениями и времена обслуживания для заявок типа $l$ являются независимыми одинаково распределенными случайными величинами с функциями распределения (CDF) $A_{l}(x)$ и $B_{l}(x)$ соответственно. Главное отличие от просто ограниченной по ресурсу модели является то, что объем требуемого ресурса представляет собой некоторую случайную величину. 

Пусть $\xi(t)$ --- количество заявок в системе в момент времени $t$, вектор $\delta(t)=(\delta_{1}(t),...,\delta_{M}(t))$ --- суммарный объем занятых ресурсов всех типов. При поступлении новой заявки с требованием $r_{i}$ в систему в момент $t_{i}$ наступает отказ и заявка теряется если отсутствует свободный сервер $\xi(t_{i}-0) = N$ или имеется нехватка свободного объема хотя бы одного из требуемых ресурсов $\delta(t_{i}-0) + r_{i} \ge R$. Иначе заявка принимается к обслуживанию и занимает объем $r_{i}$ на все время своей обработки. При завершении обслуживания заявка одновременно освобождает и сервер, и ресурсы.

При появлениии ресурсного требования у заявки систему становится крайне трудно анализировать аналитически, так как формулы стационарных вероятностей содержат многократные свертки функции распределения требований к ресурсам $F^{(k)}(x)$, что делает вычислительные алгоритмы слишком сложными. Однако, у такой СМО по прежнему сохраняется ранее озвученное свойство инвариантности к распределению времени работы заявки при Пуассоновском входящем потоке. 

Необходимость применения подобных моделей необходима, например, для описания ранее упомянутой концепции сетевых срезов в современных сетях 5G. В рамках одной физической инфраструктуры сосуществуют потоки заявок (трафика) с кардинально требованиями:
\begin{itemize}
    \item \textbf{URLLC (Ultra-Reliable Low-Latency Communication):} заявки критичны к задержкам. Они могут требовать небольшого объема радиоресурсов, но нуждаются в немедленном выделении вычислительных мощностей (например, $n$ виртуальных процессоров) для мгновенной обработки на границе сети.
    \item \textbf{eMBB (Enhanced Mobile Broadband):} заявки, требующие передачи тяжелого контента (например, 4K/8K видео). Для них критически важен вектор пропускной способности (широкая полоса радиочастот), тогда как требования к процессорному времени могут быть менее жесткими.
    \item \textbf{mMTC (Massive Machine Type Communications):} массовый поток заявок от устройств интернета вещей (IoT), где каждая отдельная заявка потребляет минимум ресурсов $r_i$, но интенсивность самого потока $\lambda$ крайне высока.
\end{itemize}

В таких условиях классическая логика ТМО с занятием только лишь "канала обслуживания" не способна адекватно описать процесс потери заявок, так как отказ для тяжелой видео-сессии (eMBB) наступит из-за исчерпания частотного ресурса, а для критического датчика (URLLC) --- из-за нехватки свободных вычислительных слотов сервера, что делает применение многомерного ресурсного вектора $\delta(t)$ строго обязательным для корректного моделирования. Тем не менее данная работа посвящена скорее чувствительности к распределению времени обслуживания и типу входящего потока, так что ресурсное требование взято за случайную величину, выраженную числом.
% Переходим от классики к реальности — вводим дополнительное ограничение по ресурсу R (оперативная память, частотный спектр и т.д.). Показываем переход к M/G/N/N/R. Описываем, как свойство инвариантности сохраняется или видоизменяется в таких системах. Можно упомянуть интенсивности, зависящие от состояния, как способ аналитического приближения. Приводим примеры применения (например, частотный ресурс в сетях 5G). Кроме того, необходимо описать связь ресурсных систем с классическими loss-system и указать, какие результаты об инвариантности уже известны для этого класса моделей.

% Блок 1: ресурсная постановка и описание состояния системы.
% Блок 2: связь с многолинейной системой Эрланга с потерями.
% Блок 3: известные результаты для ресурсных СМО.
% Блок 4: вычислительные трудности точного анализа.

% Цитировать Статьи Э.С. Сопина (включая работу 2018 года по инвариантности ресурсных систем и 2019 года со случайными сигналами), Daraseliya et al. (про 5G NR-U), Abouee-Mehrizi, Bekker.

\section{Влияние немарковских потоков на характеристики обслуживания}
\label{sec:non_markovian_service}

Классические результаты для систем с отказами и ресурсных СМО в существенной степени опираются на предположение о пуассоновском входящем потоке. Это свойство существенно упрощает аналитическое описание системы и позволяет получать стационарные характеристики в явном или рекуррентном виде. Однако современные телекоммуникационные системы от сетей с коммутацией каналов к пакетным сетям передачи данных (IP-сетям, интернету вещей, видео-стримингу) не имеют стационарности, ординарности и отсутствия последействия (строгая марковость), на которых базируются классические модели $M/M/1$ или $M/M/n$, как показывают эмпирические исследования. В таких случаях входящий поток обозначают как $G$ (General) или $GI$ (General Indepedant) для произвольных независимых интервалов прихода заявки. Именно такая постановка выбрана в работе. 

Поступления пакетов могут быть более регулярными, чем в пуассоновском случае, либо, наоборот, более неравномерными и всплесковыми. В первом случае интервалы $A$ между соседними поступлениями имеют меньшую вариативность, во втором — поток содержит чередование коротких интервалов интенсивных поступлений и длинных пауз. Для количественной оценки степени отличия таких потоков используется квадрат коэффициента вариации интервалов между поступлениями $c^2_A$. При пуассоновском потоке $c^2_A = 1$, для более плавных режимов поступления $c^2_A < 1$, для менее стабильных $c^2_A > 1$. При одинаковой средней интенсивности такие потоки могут по-разному нагружать систему и приводить к различным значениям вероятности отказа и среднего занятого ресурса. На основании значений квадрата коэффициента вариации в работе будут выбраны разные семейства распределений. 

Для описания сложных входящих потоков в теории массового обслуживания используются марковские процессы поступления (MAP), их групповые аналоги (BMAP),матрично-экспоненциальные и рациональные потоки. Эти модели позволяют учитывать фазовую структуру поступлений, корреляцию между событиями и групповой характер заявок. При этом сохраняется возможность вычислительного анализа с использованием матричных методов, но это сильно усложняет пространство состояний, особенно в ресурсных системах. Аналитические методы здесь приводит к проклятию размерности. Это обуславливает использование в работе использование имитационного моделирования для анализа стационарных характеристик и уровня чувствительности. 

% Здесь мы ломаем «идеальную» картину. Что происходит, если входящий поток перестает быть пуассоновским (M→GI) - что ломается? Описываем регулярные (Эрланг) и пачечные/нерегулярные (гиперэкспоненциальные) потоки. Показываем, что система теряет марковское свойство, а попытка решить её аналитически (через метод фиктивных фаз) приводит к «проклятию размерности» (state space explosion), особенно при наличии ресурсного ограничения

% Цитировать О марковских и рациональных потоках случайных событий (I и II), Dudin et al. (про случайную среду и фазовые распределения).

\section{Состояние вопроса и постановка задачи}
\label{sec:problem_statement}

Проведенный в предыдущих разделах обзор литературы позволяет выделить несколько ключевых этапов в развитии моделей массового обслуживания с отказами. Для базовых систем (классические модели Эрланга, описываемые в нотации Кендалла как $M/M/N/0$ и $M/G/N/N$) существует удобный аналитический аппарат. Важнейшим для данной работы результатом является теорема об инвариантности стационарных вероятностей состояний к форме распределения времени обслуживания при пуассоновском входящем потоке. Так система $M/G/N/N$ может быть удобно проанализирована аналитически, что будет проверено в ходе эксперимента. 

Современные сети связи больше не описываются только каналом, в связи с чем потребовался переход к более сложным ресурсным и мультиресурсным моделям. В таких моделях заявка требует не только прибора, но и случайного объема дополнительных ресурсов, выраженного числом или вектором. Несмотря на сохранение свойства инвариантности при пуассоновском входящем потоке, точный аналитический расчет стационарных характеристик становится вычислительно неэффективным из-за необходимости вычисления многократных сверток функций распределения требуемых ресурсов. 

Также в силу особенностей современных сетей приходится отказаться от пуассоновского потока заявок для описания моделей. Сразу теряется марковость входящего потока и инвариантность распределения времени работы. Эмпирические данные свидетельствуют, что трафик в современных сетях носит существенно немарковский характер, обладая свойствами пачечности пакетов или, напротив, повышенной регулярностью. Существующие математические аппараты для описания таких потоков (марковские процессы поступления MAP, BMAP, матрично-экспоненциальные распределения) при наложении на многомерное пространство состояний ресурсной системы приводят к экспоненциальному росту сложности вычислений, известному как «проклятие размерности».

В связи с этим исследуемая в данной работе постановка в нотации Кендалла имеет вид в $GI/G/N/N/R$ и требует перехода от строгих аналитических методов к методам имитационного моделирования\cite{sopin2018_simulation_limited_resources,sopin2019_ecms_signals}. Основной фокус сделан на чувствительности системы к семействам распределения входящего потока и времени работы, а также их взаимному влиянию; ограниченный ресурс задается числом. Поэтому отказ в обслуживании поступившей $i$-й заявки с ресурсным требованием $r_i$ может быть вызван двумя независимыми причинами\cite{naumov_total_resources_2016,naumov_samouylov_multiresource_loss_2017}:
$$ \xi(t_i-0) = N \quad \text{или} \quad \delta(t_i-0) + r_i > R. $$

Так как входящий поток рассматривается не только в пуассоновском случае даже при одинаковой средней интенсивности $\lambda$ более регулярный и более всплесковый потоки могут по-разному нагружать систему\cite{paxson_floyd_1995_poisson_failure,vishnevskii_dudin_correlated_arrivals_2017} и приводить к различным значениям вероятности отказа, пропускной способности и среднего занятого ресурса, а также вызывать разные уровни чувствительности к распределению времени работы заявок в силу отсутствия инвариантности.

Для достижения цели работы решаются следующие задачи:
\begin{enumerate}
    \item сформулировать математическую модель ресурсной системы с отказами с ограничениями по числу приборов и по суммарному ресурсу;
    \item определить набор исследуемых стационарных характеристик: вероятность отказа, пропускную способность, среднее число заявок в системе, средний занятый ресурс;
    \item задать семейства входящих потоков и распределений времени обслуживания с одинаковыми первыми моментами, но различными коэффициентами вариации;
    \item разработать дискретно-событийный алгоритм для имитационного моделирования системы;
    \item реализовать программно-вычислительный комплекс для запуска независимых репликаций экспериментов и сбора статистических данных;
    \item провести вычислительные эксперименты и проанализировать чувствительность метрик к изменению вероятностных распределений.
\end{enumerate}

% Это кульминация главы. Раз аналитика пасует перед GI/G/N/N/R, нам нужно дискретно-событийное моделирование (DES). Кратко описываем суть DES. Отмечаем проблему скорости и шума Монте-Карло, что диктует потребность в параллельных вычислениях. В конце этого подраздела формулируется строгая постановка задачи вашей диссертации: разработать программно-вычислительный комплекс для оценки чувствительности системы GI/G/N/N/R к дисперсии потоков.

% цитировать Sopin 2018 (On the Simulations of the Limited Resources Queueing Systems), репозиторий «Основы моделирования дискретных систем».
